\section{Conclusion}
\label{sec:conclusion}
This demonstration separates an attention expression, a storage schedule, a synthetic performance model, and an actual authoring artifact. The separation is the main lesson. Fewer transferred bytes, lower temporary workspace, and shorter elapsed time are different claims that require different evidence. A compact model can make their relationship visible, provided its assumptions remain attached to its results.

The numerical study shows how launch overhead can limit small-workload improvements, how arithmetic can become the active constraint after traffic is reduced, and how shape and resource assumptions change a ratio. These are conditional consequences of the specified model. They are not measured accelerator results and should not be cited as such. The project supplies a reproducible starting point for further analysis, including the source data, figures, BibTeX records, and compiled PDF; the companion generator is included in the project tools directory.

A substantive next step would be to replace the abstract schedules with validated implementations and collect measurements under a documented protocol. Until then, the artifact is useful as a reasoning aid and a demonstration of transparent scientific authoring. Its value comes from making the boundary between illustration and evidence easy to inspect.
